% ---
% toc: true
% ---
\documentclass[11pt]{scrartcl}

\usepackage[magyar]{babel}
\usepackage[stixtwo]{fontsetup}
\usepackage{amsmath}
\usepackage{amsthm}
\usepackage{booktabs}
\usepackage[margin=1in]{geometry}
\usepackage{tikz}
\usepackage{pgfplots}

\usetikzlibrary{arrows.meta}
\pgfplotsset{compat=1.18}
\usepgfplotslibrary{groupplots}
\usepackage{luacode}

% Feliratméretek. A weboldalon (lwarp-pal konvertálva) az ábrák kicsinyítve
% jelennek meg, ott nagyobb feliratok kellenek, mint a PDF-ben.
\makeatletter
\@ifpackageloaded{lwarp}{%
	\newcommand{\figfont}{\Large}%
	\newcommand{\figtickfont}{\Large}%
}{%
	\newcommand{\figfont}{\normalsize}%
	\newcommand{\figtickfont}{\small}%
}
% Figures with meaningful colours (heat maps) are wrapped in
% \begin{BlockClass}{keep-colors}: the website does not invert them in the dark
% theme. Without lwarp (in the PDF) the environment does nothing.
\@ifpackageloaded{lwarp}{}{\newenvironment{BlockClass}[2][]{}{}}
\makeatother

\setlength{\parindent}{0pt}
\setlength{\parskip}{0.66em}

\DeclareMathOperator{\arcosh}{arcosh}
\DeclareMathOperator{\Real}{Re}
\DeclareMathOperator{\Imag}{Im}

% Egy szín minden függvénynek, az összes ábrán (mint a másik két jegyzetben).
\definecolor{sincolor}{HTML}{1F77B4}
\definecolor{coscolor}{HTML}{D62728}
\definecolor{tancolor}{HTML}{2CA02C}
\definecolor{arccolor}{HTML}{FF7F0E}
\definecolor{invcolor}{HTML}{6A3D9A}

% The surfaces of the 3D figures. \complexpanel{sin}{re} starts the next plot
% of a groupplot with the real part of sin z (parts: re, im, abs) and draws
% its slice along the real axis (along the imaginary axis for the hyperbolic
% functions) in red. A point of the slice is drawn only if the surface does
% not hide it: the ray from it towards the viewer is followed, and if the
% surface crosses the ray, the point is hidden (a depth test, as a GPU's
% z-buffer does, but exact for the height field). pgfplots projects the axis
% box as a unit cube, so the test runs in those normalised coordinates.
\begin{luacode*}
local VIEW_AZ, VIEW_EL = math.rad(25), math.rad(35)
local LONG, SHORT = 2 * math.pi, 1.6

local function sinh(t) return (math.exp(t) - math.exp(-t)) / 2 end
local function cosh(t) return (math.exp(t) + math.exp(-t)) / 2 end

-- f(x + iy) as its real and imaginary part.
local functions = {
  sin = function(x, y) return math.sin(x) * cosh(y), math.cos(x) * sinh(y) end,
  cos = function(x, y) return math.cos(x) * cosh(y), -math.sin(x) * sinh(y) end,
  tan = function(x, y)
    local d = math.cos(2 * x) + cosh(2 * y)
    return math.sin(2 * x) / d, sinh(2 * y) / d
  end,
  sinh = function(x, y) return sinh(x) * math.cos(y), cosh(x) * math.sin(y) end,
  cosh = function(x, y) return cosh(x) * math.cos(y), sinh(x) * math.sin(y) end,
  tanh = function(x, y)
    local d = cosh(2 * x) + math.cos(2 * y)
    return sinh(2 * x) / d, math.sin(2 * y) / d
  end,
}

local titles = { re = "$\\Real \\%s z$", im = "$\\Imag \\%s z$", abs = "$|\\%s z|$" }

local function fmt(x, y, z) return string.format("(%.4f,%.4f,%.4f)", x, y, z) end

function complexpanel(name, part)
  local hyperbolic = name:sub(-1) == "h"
  -- The tangents are cut at 3, near their poles.
  local cut = (name == "tan" or name == "tanh") and 3 or 2.6
  local box = hyperbolic
      and { xmin = -SHORT, xmax = SHORT, ymin = -LONG, ymax = LONG }
      or { xmin = -LONG, xmax = LONG, ymin = -SHORT, ymax = SHORT }
  box.zmin, box.zmax = part == "abs" and 0 or -cut, cut

  local function value(x, y)
    local re, im = functions[name](x, y)
    if part == "re" then return re end
    if part == "im" then return im end
    return math.sqrt(re * re + im * im)
  end
  local function surface(x, y)
    local z = value(x, y)
    if z ~= z then return box.zmax end
    return math.max(box.zmin, math.min(box.zmax, z))
  end

  local ztick = part == "abs" and (cut == 3 and "{0,1,2,3}" or "{0,1,2}")
      or (cut == 3 and "{-3,0,3}" or "{-2,0,2}")
  tex.sprint(string.format(
    "\\nextgroupplot[%s, title={%s}, xmin=%f, xmax=%f, ymin=%f, ymax=%f, zmin=%f, zmax=%f, ztick=%s]",
    hyperbolic and "hyperbolic surface" or "trig surface",
    string.format(titles[part], name),
    box.xmin, box.xmax, box.ymin, box.ymax, box.zmin, box.zmax, ztick))

  -- The surface; 100 steps along the long side, so that no grid point falls
  -- on a pole of a tangent.
  local nx, ny = 101, 41
  if hyperbolic then nx, ny = 41, 101 end
  local points = {}
  for j = 0, ny - 1 do
    local y = box.ymin + (box.ymax - box.ymin) * j / (ny - 1)
    for i = 0, nx - 1 do
      local x = box.xmin + (box.xmax - box.xmin) * i / (nx - 1)
      points[#points + 1] = fmt(x, y, surface(x, y))
    end
  end
  tex.sprint(string.format(
    "\\addplot3[surf, shader=interp, mesh/cols=%d] coordinates {%s};",
    nx, table.concat(points, " ")))

  -- The depth test, in the unit cube of the axis box.
  local view = {
    math.sin(VIEW_AZ) * math.cos(VIEW_EL),
    -math.cos(VIEW_AZ) * math.cos(VIEW_EL),
    math.sin(VIEW_EL),
  }
  -- The point lies on the surface, so the ray starts on one side of it (above
  -- it, or below it where we see the surface from underneath); the point is
  -- hidden if the ray crosses the surface later.
  local function hidden(x, y, z)
    local u = (x - box.xmin) / (box.xmax - box.xmin)
    local v = (y - box.ymin) / (box.ymax - box.ymin)
    local w = (z - box.zmin) / (box.zmax - box.zmin)
    local side
    for t = 0.002, 2, 0.002 do
      local tu, tv, tw = u + t * view[1], v + t * view[2], w + t * view[3]
      if tu < 0 or tu > 1 or tv < 0 or tv > 1 then return false end
      local sz = surface(box.xmin + tu * (box.xmax - box.xmin),
        box.ymin + tv * (box.ymax - box.ymin))
      local gap = (sz - box.zmin) / (box.zmax - box.zmin) - tw
      if side == nil then
        side = gap > 0 and 1 or -1
      elseif gap * side < -1e-4 then
        return true
      end
    end
    return false
  end

  -- The slice, split into its visible runs; it also breaks where it leaves
  -- the box (at the poles of the tangents).
  local runs, run = {}, {}
  local n = 401
  for k = 0, n - 1 do
    local t = -LONG + 2 * LONG * k / (n - 1)
    local x, y = t, 0
    if hyperbolic then x, y = 0, t end
    local z = value(x, y)
    if z ~= z or z < box.zmin or z > box.zmax or hidden(x, y, z) then
      if #run > 1 then runs[#runs + 1] = run end
      run = {}
    else
      run[#run + 1] = fmt(x, y, z)
    end
  end
  if #run > 1 then runs[#runs + 1] = run end
  for _, r in ipairs(runs) do
    tex.sprint("\\addplot3[axis curve] coordinates {" .. table.concat(r, " ") .. "};")
  end
end
\end{luacode*}
\newcommand{\complexpanel}[2]{\directlua{complexpanel("#1", "#2")}}

\theoremstyle{definition}
\newtheorem{definition}{Definíció}
\newtheorem{theorem}[definition]{Tétel}
\newtheorem{example}[definition]{Példa}
\theoremstyle{remark}
\newtheorem*{remark}{Megjegyzés}

% Közös stílus a függvénygrafikonokhoz.
\pgfplotsset{
	% |f(z)| over the complex plane, as a heat map seen from above.
	modulus/.style={
			view={0}{90},
			width=0.42\linewidth,
			height=0.42\linewidth,
			xmin=-6.2832, xmax=6.2832, ymin=-6.2832, ymax=6.2832,
			domain=-6.2832:6.2832,
			y domain=-6.2832:6.2832,
			% 60 steps: no grid point falls on a pole of tan or tanh.
			samples=61,
			enlargelimits=false,
			axis on top,
			trig format plots=rad,
			colormap/viridis,
			% A logarithmic colour scale: the plots show ln(1 + |f(z)|), so that
			% both the zeros and the fast growth are visible; the colour bar is
			% labelled with |f(z)|.
			point meta min=0,
			point meta max=5.5947,
			xtick={-6.2832, -3.1416, 0, 3.1416, 6.2832},
			xticklabels={$-2\pi$, $-\pi$, $0$, $\pi$, $2\pi$},
			ytick={-6.2832, -3.1416, 0, 3.1416, 6.2832},
			yticklabels={$-2\pi$, $-\pi$, $0$, $\pi$, $2\pi$},
			xlabel={$\Real z$},
			ylabel={$\Imag z$},
			ticklabel style={font=\figtickfont},
			label style={font=\figfont},
			title style={font=\figfont},
		},
	% A part of f(z) as a surface over a strip of the complex plane; the
	% surfaces are computed by \complexpanel below. The view must match
	% VIEW_AZ and VIEW_EL there.
	surface/.style={
			width=0.5\linewidth,
			height=0.39\linewidth,
			view={25}{35},
			z buffer=sort,
			colormap/viridis,
			ticklabel style={font=\figtickfont},
			label style={font=\figfont},
			title style={font=\figfont},
			xlabel={$\Real z$},
			ylabel={$\Imag z$},
		},
	trig surface/.style={
			xtick={-6.2832, -3.1416, 0, 3.1416, 6.2832},
			xticklabels={$-2\pi$, $-\pi$, $0$, $\pi$, $2\pi$},
			ytick={-1, 0, 1},
		},
	hyperbolic surface/.style={
			xtick={-1, 0, 1},
			ytick={-6.2832, -3.1416, 0, 3.1416, 6.2832},
			yticklabels={$-2\pi$, $-\pi$, $0$, $\pi$, $2\pi$},
		},
	% The slice of a surface along the real axis (trigonometric) or the
	% imaginary axis (hyperbolic).
	axis curve/.style={red!85!black, line width=1.6pt},
	zeros/.style={only marks, mark=*, mark size=2.2pt, white,
			mark options={draw=black, line width=0.6pt}},
	poles/.style={only marks, mark=diamond*, mark size=3.2pt, black,
			mark options={draw=white, line width=0.6pt}},
	modulus colorbar/.style={
			colorbar,
			colorbar style={
					ytick={0, 1.0986, 2.3979, 3.9318, 5.5255},
					yticklabels={$0$, $2$, $10$, $50$, $250$},
					yticklabel style={font=\figtickfont},
				},
		},
	% The tangents stay near 1 away from their poles, so they get a finer
	% scale of their own, up to 10.
	tangent scale/.style={
			point meta max=2.3979,
			colorbar style={
					ytick={0, 0.6931, 1.0986, 1.6094, 2.3979},
					yticklabels={$0$, $1$, $2$, $4$, $10$},
				},
		},
	functionplot/.style={
			axis lines=middle,
			axis line style={-{Stealth[scale=1.1]}, gray!70!black, line width=1pt},
			xlabel={$x$},
			ylabel={$y$},
			xlabel style={anchor=west},
			ylabel style={anchor=south},
			tick style={gray!70!black},
			ticklabel style={font=\figtickfont, fill=white, inner sep=1pt},
			label style={font=\figfont},
			trig format plots=rad,
			% A tengelyek és a feliratok (fehér háttéren) a görbék fölött.
			axis on top,
			samples=200,
			clip=false,
		},
}
\tikzset{
	% A TikZ alapértelmezésénél vastagabb vonalak: a weboldalon az ábrák
	% kicsinyítve jelennek meg.
	every picture/.append style={line width=0.9pt, font=\figfont},
	% Ha egy görbe kifut az ábrából, ott nyílhegyben végződik; ha valóban
	% véget ér, ott pontban (end dot).
	tip/.tip={Stealth[length=8pt, width=7pt]},
	continues/.style={tip-tip},
	continues right/.style={-tip},
	end dot/.style={only marks, mark=*, mark size=2.8pt},
	graph/.style={line width=2.2pt, line cap=round},
	faded/.style={line width=1.6pt, gray!55},
	mirror/.style={densely dashed, line width=1pt, gray!75},
}

\title{A trigonometrikus és a hiperbolikus függvények kapcsolata}
\subtitle{Analízis A1 -- segédanyag}
\author{Sándor Tibor}
\date{}

\begin{document}
\maketitle

Az előző két jegyzetben a trigonometrikus és a hiperbolikus függvények
feltűnően hasonlóan viselkedtek: az egységkör és az egységhiperbola, a két
Pitagorasz-tétel, a párosság és a páratlanság mind párban jártak. Ez nem
véletlen. A komplex számok körében a két függvénycsalád lényegében ugyanaz,
csak a képzetes tengely mentén elforgatva:
\[
	\cos x = \cosh(ix), \qquad i \sin x = \sinh(ix)\text.
\]
Ebben a jegyzetben ezt az összefüggést mutatjuk meg, és megnézzük, mire jó.

\section{Az Euler-formula}

A komplex számok exponenciális alakjából ismert az \emph{Euler-formula}:
minden $x \in \mathbb{R}$ esetén
\[
	e^{ix} = \cos x + i \sin x\text.
\]
Mivel a koszinusz páros, a szinusz pedig páratlan, $x$ helyére $-x$-et írva
\[
	e^{-ix} = \cos x - i \sin x\text.
\]
A két egyenlőséget összeadva, illetve kivonva a koszinusz és a szinusz az
exponenciális függvénnyel is kifejezhető (\ref{fig:euler}.~ábra):
\[
	\cos x = \frac{e^{ix} + e^{-ix}}{2}, \qquad
	\sin x = \frac{e^{ix} - e^{-ix}}{2i}\text.
\]
Vessük össze ezt a hiperbolikus függvények exponenciális alakjával:
\[
	\cosh x = \frac{e^{x} + e^{-x}}{2}, \qquad
	\sinh x = \frac{e^{x} - e^{-x}}{2}\text.
\]
A különbség csak annyi, hogy a trigonometrikus képletekben $e^{x}$ helyett
$e^{ix}$ áll, a szinusz nevezőjében pedig $2$ helyett $2i$.

\begin{figure}[ht]
	\centering
	\begin{tikzpicture}[scale=2.4, line cap=round]
		\def\x{1.1}
		\coordinate (O) at (0, 0);
		\coordinate (E) at ({deg(\x)}:1);
		\coordinate (F) at ({-deg(\x)}:1);
		\coordinate (C) at ({cos(deg(\x))}, 0);

		\draw[-{Stealth}, gray!70!black] (-1.3, 0) -- (1.5, 0)
		node[below] {$\Real$};
		\draw[-{Stealth}, gray!70!black] (0, -1.3) -- (0, 1.35)
		node[left] {$\Imag$};
		\draw[line width=1.5pt] (O) circle (1);
		\draw (O) -- (E);
		\draw (O) -- (F);

		\draw[->, arccolor] (0.25, 0) arc (0:{deg(\x)}:0.25);
		\node[arccolor] at ({deg(\x) / 2}:0.4) {$x$};

		\draw[mirror, sincolor] (C) -- (F);
		\draw[line width=2.8pt, coscolor] (O) -- (C);
		\draw[line width=2.8pt, sincolor] (C) -- (E);
		\node[sincolor, right] at ({cos(deg(\x))}, {sin(deg(\x)) / 2})
		{$\sin x$};
		\node[coscolor, below] at ({cos(deg(\x)) / 2}, 0) {$\cos x$};

		\fill (E) circle (1.4pt) node[above right] {$e^{ix}$};
		\fill (F) circle (1.4pt) node[below right] {$e^{-ix}$};
		\fill (C) circle (1.4pt);
	\end{tikzpicture}
	\caption{Az $e^{ix}$ és az $e^{-ix}$ pont egymás tükörképe a valós
		tengelyre. Összegük valós: $2\cos x$ (a piros szakasz $\cos x$),
		különbségük tisztán képzetes: $2i \sin x$ (a kék szakasz $\sin x$).}
	\label{fig:euler}
\end{figure}

\section{Kiterjesztés a komplex számokra}

A komplex kitevőjű hatványt a komplex számok exponenciális alakja már
értelmezi: ha $z = a + bi$, ahol $a, b \in \mathbb{R}$, akkor
\[
	e^{z} = e^{a} (\cos b + i \sin b)\text.
\]
Az előző fejezet képletei így tetszőleges komplex számra is értelmesek, és
ezekkel terjesztjük ki a függvényeket.

\begin{definition}
	Minden $z \in \mathbb{C}$ esetén
	\[
		\cos z := \frac{e^{iz} + e^{-iz}}{2}, \qquad
		\sin z := \frac{e^{iz} - e^{-iz}}{2i}, \qquad
		\cosh z := \frac{e^{z} + e^{-z}}{2}, \qquad
		\sinh z := \frac{e^{z} - e^{-z}}{2}\text,
	\]
	továbbá $\tan z := \sin z/\cos z$ és
	$\tanh z := \sinh z/\cosh z$, ahol a nevező nem nulla.
\end{definition}

Valós $z$ esetén ezek a megszokott függvények: a trigonometrikusak az
Euler-formula, a hiperbolikusak az exponenciális alakjuk miatt.

\section{A kapcsolat}

\begin{theorem}
	Minden $z \in \mathbb{C}$ esetén
	\[
		\cosh(iz) = \cos z, \qquad \sinh(iz) = i \sin z,
	\]
	\[
		\cos(iz) = \cosh z, \qquad \sin(iz) = i \sinh z,
	\]
	és ahol a nevező nem nulla,
	\[
		\tanh(iz) = i \tan z, \qquad \tan(iz) = i \tanh z\text.
	\]
\end{theorem}

\begin{proof}
	Helyettesítsünk be a definíciókba. Az első kettő:
	\[
		\cosh(iz) = \frac{e^{iz} + e^{-iz}}{2} = \cos z, \qquad
		\sinh(iz) = \frac{e^{iz} - e^{-iz}}{2}
		= i \cdot \frac{e^{iz} - e^{-iz}}{2i} = i \sin z\text.
	\]
	A második kettőnél $i \cdot iz = -z$:
	\[
		\cos(iz) = \frac{e^{-z} + e^{z}}{2} = \cosh z, \qquad
		\sin(iz) = \frac{e^{-z} - e^{z}}{2i}
		= i \cdot \frac{e^{z} - e^{-z}}{2} = i \sinh z\text,
	\]
	ahol felhasználtuk, hogy $1/i = -i$. A tangensekre ezekből
	\[
		\tanh(iz) = \frac{\sinh(iz)}{\cosh(iz)} = \frac{i \sin z}{\cos z}
		= i \tan z, \qquad
		\tan(iz) = \frac{\sin(iz)}{\cos(iz)} = \frac{i \sinh z}{\cosh z}
		= i \tanh z\text.
	\]
\end{proof}

Szavakban: a koszinusz a képzetes tengely mentén koszinusz hiperbolikusz, a
szinusz pedig ($i$-szeres) szinusz hiperbolikusz. Valós $y$ esetén például
$\cos(iy) = \cosh y$, ami $y \to \pm\infty$ esetén végtelenhez tart
(\ref{fig:cos-imaginary}.~ábra).

\begin{figure}[ht]
	\centering
	\begin{tikzpicture}
		\begin{axis}[
				functionplot,
				width=\linewidth,
				height=7cm,
				xmin=-6.9, xmax=7.2,
				ymin=-1.5, ymax=4.6,
				xtick={-6, -4, -2, 2, 4, 6},
				ytick={-1, 1, 2, 3, 4},
			]
			\addplot[graph, continues, coscolor, domain=-6.6:6.6] {cos(x)}
			node[pos=0.97, below] {$\cos x$};
			\addplot[graph, continues, invcolor, domain=-2.1:2.1] {cosh(x)}
			node[pos=0.95, right] {$\cos(ix) = \cosh x$};
			\addplot[end dot, black] coordinates {(0, 1)};
		\end{axis}
	\end{tikzpicture}
	\caption{A valós tengely mentén a koszinusz $-1$ és $1$ között hullámzik
		(piros), a képzetes tengely mentén viszont $\cos(ix) = \cosh x$
		(lila), ami nem korlátos. Mindkettő páros, és értéke a $0$-ban $1$.}
	\label{fig:cos-imaginary}
\end{figure}

\begin{remark}
	Valós számokra $|\cos x| \le 1$ és $|\sin x| \le 1$, a komplex síkon
	azonban ez már nem igaz: például $\cos i = \cosh 1 \approx 1{,}54$, és
	$|\sin(iy)| = |\sinh y|$ tetszőlegesen nagy lehet.
\end{remark}

\section{Azonosságok átvitele}

A kapcsolat haszna, hogy a trigonometrikus azonosságokból számolás nélkül
kapunk hiperbolikusakat. A trigonometrikus addíciós tételek, például
\[
	\sin(z + w) = \sin z \cos w + \cos z \sin w, \qquad
	\cos(z + w) = \cos z \cos w - \sin z \sin w\text,
\]
az exponenciális alakból ugyanúgy beláthatók komplex $z$ és $w$ esetén is,
mint valós számokra.

\begin{example}
	Vezessük le a trigonometrikus Pitagorasz-tételből a hiperbolikusat!

	Írjunk $x$ helyére $ix$-et a $\cos^2 x + \sin^2 x = 1$ azonosságba:
	\[
		1 = \cos^2(ix) + \sin^2(ix)
		= \cosh^2 x + (i \sinh x)^2
		= \cosh^2 x - \sinh^2 x\text.
	\]
\end{example}

\begin{example}
	Vezessük le a hiperbolikus addíciós tételeket!

	Írjunk $x$ helyére $ix$-et, $y$ helyére $iy$-t:
	\[
		\begin{aligned}
			\cosh(x + y) & = \cos(ix + iy)
			= \cos(ix) \cos(iy) - \sin(ix) \sin(iy)                         \\
			             & = \cosh x \cosh y - (i \sinh x)(i \sinh y)\text,
		\end{aligned}
	\]
	vagyis
	\[
		\cosh(x + y) = \cosh x \cosh y + \sinh x \sinh y\text.
	\]
	Hasonlóan, $\sinh(x + y) = -i \sin(ix + iy)$ miatt
	\[
		\begin{aligned}
			\sinh(x + y) & = -i \bigl(\sin(ix) \cos(iy) + \cos(ix) \sin(iy)\bigr)       \\
			             & = -i \bigl(i \sinh x \cosh y + i \cosh x \sinh y\bigr)\text,
		\end{aligned}
	\]
	vagyis
	\[
		\sinh(x + y) = \sinh x \cosh y + \cosh x \sinh y\text.
	\]
\end{example}

A példákban mindig ugyanaz történt: a $\cos$ helyére $\cosh$, a $\sin$ helyére
$\sinh$ került, és ahol két szinusz szorzata állt, ott $i \cdot i = -1$ miatt
előjelet váltott a tag. Ez az \emph{Osborn-szabály}:

\begin{center}
	\renewcommand{\arraystretch}{1.7}
	\begin{tabular}{@{}ll@{}}
		\toprule
		Trigonometrikus                               & Hiperbolikus \\
		\midrule
		$\cos^2 x + \sin^2 x = 1$                     &
		$\cosh^2 x - \sinh^2 x = 1$                                  \\
		$\sin 2x = 2 \sin x \cos x$                   &
		$\sinh 2x = 2 \sinh x \cosh x$                               \\
		$\cos 2x = \cos^2 x - \sin^2 x$               &
		$\cosh 2x = \cosh^2 x + \sinh^2 x$                           \\
		$\sin(x + y) = \sin x \cos y + \cos x \sin y$ &
		$\sinh(x + y) = \sinh x \cosh y + \cosh x \sinh y$           \\
		$\cos(x + y) = \cos x \cos y - \sin x \sin y$ &
		$\cosh(x + y) = \cosh x \cosh y + \sinh x \sinh y$           \\
		$1 + \tan^2 x = \dfrac{1}{\cos^2 x}$                   &
		$1 - \tanh^2 x = \dfrac{1}{\cosh^2 x}$                                \\
		\bottomrule
	\end{tabular}
\end{center}

A $\tan^2 x = \sin^2 x/\cos^2 x$ is két szinusz szorzatát tartalmazza,
ezért vált előjelet a tangenses azonosság.

\section{Valós és képzetes rész}

Az addíciós tételekből a komplex szinusz és koszinusz valós és képzetes része
is kiszámolható.

\begin{theorem}
	Minden $x, y \in \mathbb{R}$ esetén
	\[
		\sin(x + iy) = \sin x \cosh y + i \cos x \sinh y, \qquad
		\cos(x + iy) = \cos x \cosh y - i \sin x \sinh y\text.
	\]
\end{theorem}

\begin{proof}
	Az addíciós tételekből, $\cos(iy) = \cosh y$ és $\sin(iy) = i \sinh y$
	felhasználásával:
	\[
		\sin(x + iy) = \sin x \cos(iy) + \cos x \sin(iy)
		= \sin x \cosh y + i \cos x \sinh y\text,
	\]
	\[
		\cos(x + iy) = \cos x \cos(iy) - \sin x \sin(iy)
		= \cos x \cosh y - i \sin x \sinh y\text.
	\]
\end{proof}

\begin{example}
	Oldjuk meg a $\cos z = 2$ egyenletet a komplex számok körében!

	Legyen $z = x + iy$, ahol $x, y \in \mathbb{R}$. Az előző tétel szerint
	\[
		\cos x \cosh y - i \sin x \sinh y = 2\text,
	\]
	vagyis a valós és a képzetes részeket összevetve
	\[
		\cos x \cosh y = 2 \qquad\text{és}\qquad \sin x \sinh y = 0\text.
	\]
	Ha $\sinh y = 0$, akkor $y = 0$, és $\cos x = 2$ lehetetlen. Így
	$\sin x = 0$, azaz $x = k\pi$, és ekkor $\cos x = \pm 1$. Mivel
	$\cosh y \ge 1 > 0$, csak $\cos x = 1$ jöhet szóba, vagyis $x = 2k\pi$ és
	$\cosh y = 2$. Ennek megoldásai
	$y = \pm \arcosh 2 = \pm \ln\left(2 + \sqrt{3}\right)$. A megoldás tehát:
	\[
		z = 2k\pi \pm i \ln\left(2 + \sqrt{3}\right)
		\qquad (k \in \mathbb{Z})\text.
	\]
\end{example}

\begin{example}
	Hol nulla a komplex szinusz?

	Legyen $z = x + iy$. A $\sin z = 0$ egyenlet valós és képzetes része:
	\[
		\sin x \cosh y = 0 \qquad\text{és}\qquad \cos x \sinh y = 0\text.
	\]
	Mivel $\cosh y \ge 1$, az elsőből $\sin x = 0$, azaz $x = k\pi$. Ekkor
	$\cos x = \pm 1$, így a másodikból $\sinh y = 0$, azaz $y = 0$. A komplex
	szinusznak tehát nincs új zérushelye: ezek pontosan a valós
	$z = k\pi$ $(k \in \mathbb{Z})$ számok.
\end{example}

\section{A komplex függvények ábrázolása}

Egy $f\colon \mathbb{C} \to \mathbb{C}$ függvény grafikonja négydimenziós
volna: két dimenzió kell a változónak és kettő az értéknek, így közvetlenül
nem rajzolható le. Szemléletes megoldás, ha a komplex sík minden $z$ pontját
az $|f(z)|$ abszolút érték szerint színezzük, vagy $|f(z)|$-t felületként a
sík fölé rajzoljuk. Az abszolút értékek az előző fejezet tételéből
kiszámolhatók.

\begin{theorem}
	Minden $x, y \in \mathbb{R}$ esetén
	\[
		|\sin(x + iy)|^2 = \sin^2 x + \sinh^2 y, \qquad
		|\cos(x + iy)|^2 = \cos^2 x + \sinh^2 y,
	\]
	\[
		|\sinh(x + iy)|^2 = \sinh^2 x + \sin^2 y, \qquad
		|\cosh(x + iy)|^2 = \sinh^2 x + \cos^2 y\text,
	\]
	és ahol a nevező nem nulla,
	\[
		|\tan(x + iy)|^2 = \frac{\sin^2 x + \sinh^2 y}{\cos^2 x + \sinh^2 y},
		\qquad
		|\tanh(x + iy)|^2 = \frac{\sinh^2 x + \sin^2 y}{\sinh^2 x + \cos^2 y}\text.
	\]
\end{theorem}

\begin{proof}
	Az előző tétel szerint a valós és a képzetes rész négyzetösszege
	\[
		\begin{aligned}
			|\sin(x + iy)|^2 & = \sin^2 x \cosh^2 y + \cos^2 x \sinh^2 y
			= \sin^2 x \left(1 + \sinh^2 y\right) + \cos^2 x \sinh^2 y \\
			                 & = \sin^2 x + \sinh^2 y\text,
		\end{aligned}
	\]
	és ugyanígy $|\cos(x + iy)|^2 = \cos^2 x + \sinh^2 y$. A hiperbolikus
	függvényekre $\sinh z = -i \sin(iz)$ és $\cosh z = \cos(iz)$, ahol
	$i(x + iy) = -y + ix$, így
	\[
		|\sinh(x + iy)| = |\sin(-y + ix)|, \qquad
		|\cosh(x + iy)| = |\cos(-y + ix)|\text,
	\]
	és ezekre az előző két képlet adja az állítást. A tangensek abszolút
	értéke pedig a számláló és a nevező abszolút értékének hányadosa.
\end{proof}

\begin{figure}[ht]
	\centering
	\begin{BlockClass}{keep-colors}
	\begin{tikzpicture}
		\begin{groupplot}[
				modulus,
				group style={
						group size=2 by 3,
						horizontal sep=1.8cm,
						vertical sep=2.2cm,
					},
			]
			\nextgroupplot[title={$|\sin z|$}]
			\addplot3[surf, shader=interp] {ln(1 + sqrt(sin(x)^2 + sinh(y)^2))};
			\addplot3[zeros] coordinates {(-2*pi, 0, 0) (-pi, 0, 0) (0, 0, 0)
					(pi, 0, 0) (2*pi, 0, 0)};
			\nextgroupplot[title={$|\sinh z|$}, modulus colorbar]
			\addplot3[surf, shader=interp] {ln(1 + sqrt(sinh(x)^2 + sin(y)^2))};
			\addplot3[zeros] coordinates {(0, -2*pi, 0) (0, -pi, 0) (0, 0, 0)
					(0, pi, 0) (0, 2*pi, 0)};
			\nextgroupplot[title={$|\cos z|$}]
			\addplot3[surf, shader=interp] {ln(1 + sqrt(cos(x)^2 + sinh(y)^2))};
			\addplot3[zeros] coordinates {(-3*pi/2, 0, 0) (-pi/2, 0, 0)
					(pi/2, 0, 0) (3*pi/2, 0, 0)};
			\nextgroupplot[title={$|\cosh z|$}, modulus colorbar]
			\addplot3[surf, shader=interp] {ln(1 + sqrt(sinh(x)^2 + cos(y)^2))};
			\addplot3[zeros] coordinates {(0, -3*pi/2, 0) (0, -pi/2, 0)
					(0, pi/2, 0) (0, 3*pi/2, 0)};
			\nextgroupplot[title={$|\tan z|$}, tangent scale]
			\addplot3[surf, shader=interp]
			{ln(1 + sqrt((sin(x)^2 + sinh(y)^2)/(cos(x)^2 + sinh(y)^2)))};
			\addplot3[zeros] coordinates {(-2*pi, 0, 0) (-pi, 0, 0) (0, 0, 0)
					(pi, 0, 0) (2*pi, 0, 0)};
			\addplot3[poles] coordinates {(-3*pi/2, 0, 0) (-pi/2, 0, 0)
					(pi/2, 0, 0) (3*pi/2, 0, 0)};
			\nextgroupplot[title={$|\tanh z|$}, modulus colorbar, tangent scale]
			\addplot3[surf, shader=interp]
			{ln(1 + sqrt((sinh(x)^2 + sin(y)^2)/(sinh(x)^2 + cos(y)^2)))};
			\addplot3[zeros] coordinates {(0, -2*pi, 0) (0, -pi, 0) (0, 0, 0)
					(0, pi, 0) (0, 2*pi, 0)};
			\addplot3[poles] coordinates {(0, -3*pi/2, 0) (0, -pi/2, 0)
					(0, pi/2, 0) (0, 3*pi/2, 0)};
		\end{groupplot}
	\end{tikzpicture}
	\end{BlockClass}
	\caption{A hat függvény abszolút értéke a komplex síkon, logaritmikus
		színskálán (sötét: kicsi, világos: nagy). A fehér pontok a
		zérushelyek, a fekete rombuszok a tangensek pólusai (ahol a nevező
		nulla). Az alsó sor színskálája finomabb, mert a tangensek a
		pólusaiktól távol $1$ közelében maradnak. A hiperbolikus függvények képe (jobbra) a
		trigonometrikusaké (balra) $90^\circ$-kal elforgatva.}
	\label{fig:moduli}
\end{figure}

Az abszolút érték mellett a valós és a képzetes rész is felületként
ábrázolható (\ref{fig:sine-surfaces}--\ref{fig:tangent-surfaces}.~ábra). A
trigonometrikus függvényekét az előző fejezet tétele adja, a hiperbolikusakét
pedig ugyanígy az addíciós tételek:
\[
	\sinh(x + iy) = \sinh x \cos y + i \cosh x \sin y, \qquad
	\cosh(x + iy) = \cosh x \cos y + i \sinh x \sin y\text.
\]
A tangensekre a tört bővítésével (a számlálót és a nevezőt is
$\cos(x - iy)$-vel, illetve $\cosh(x - iy)$-vel szorozva)
\[
	\tan(x + iy) = \frac{\sin 2x + i \sinh 2y}{\cos 2x + \cosh 2y}, \qquad
	\tanh(x + iy) = \frac{\sinh 2x + i \sin 2y}{\cosh 2x + \cos 2y}\text.
\]

\begin{figure}[p]
	\centering
	\begin{BlockClass}{keep-colors}
	\begin{tikzpicture}
		\begin{groupplot}[
				surface,
				group style={
						group size=2 by 3,
						horizontal sep=1.2cm,
						vertical sep=2.1cm,
					},
			]
			\complexpanel{sin}{re}\complexpanel{sinh}{re}
			\complexpanel{sin}{im}\complexpanel{sinh}{im}
			\complexpanel{sin}{abs}\complexpanel{sinh}{abs}
		\end{groupplot}
	\end{tikzpicture}
	\end{BlockClass}
	\caption{A $\sin z$ (balra) és a $\sinh z$ (jobbra) valós része, képzetes
		része és abszolút értéke. A piros görbe a felület metszete a valós, illetve a képzetes tengely
		fölött; ahol a felület eltakarja, ott nem látszik. A valós tengelyen
		$\sin x$ valós, a képzetes tengelyen $\sinh(iy) = i \sin y$ tisztán
		képzetes.}
	\label{fig:sine-surfaces}
\end{figure}

\begin{figure}[p]
	\centering
	\begin{BlockClass}{keep-colors}
	\begin{tikzpicture}
		\begin{groupplot}[
				surface,
				group style={
						group size=2 by 3,
						horizontal sep=1.2cm,
						vertical sep=2.1cm,
					},
			]
			\complexpanel{cos}{re}\complexpanel{cosh}{re}
			\complexpanel{cos}{im}\complexpanel{cosh}{im}
			\complexpanel{cos}{abs}\complexpanel{cosh}{abs}
		\end{groupplot}
	\end{tikzpicture}
	\end{BlockClass}
	\caption{A $\cos z$ (balra) és a $\cosh z$ (jobbra) valós része, képzetes
		része és abszolút értéke. A piros görbe a felület metszete a valós, illetve a képzetes tengely
		fölött; ahol a felület eltakarja, ott nem látszik. Mindkét metszet valós:
		$\cos x$, illetve $\cosh(iy) = \cos y$.}
	\label{fig:cosine-surfaces}
\end{figure}

\begin{figure}[p]
	\centering
	\begin{BlockClass}{keep-colors}
	\begin{tikzpicture}
		\begin{groupplot}[
				surface,
				group style={
						group size=2 by 3,
						horizontal sep=1.2cm,
						vertical sep=2.1cm,
					},
			]
			\complexpanel{tan}{re}\complexpanel{tanh}{re}
			\complexpanel{tan}{im}\complexpanel{tanh}{im}
			\complexpanel{tan}{abs}\complexpanel{tanh}{abs}
		\end{groupplot}
	\end{tikzpicture}
	\end{BlockClass}
	\caption{A $\tan z$ (balra) és a $\tanh z$ (jobbra) valós része, képzetes
		része és abszolút értéke, a pólusok közelében $\pm 3$-nál
		levágva. A piros görbe a felület metszete a valós, illetve a képzetes tengely
		fölött; ahol a felület eltakarja, ott nem látszik. A valós tengelyen $\tan x$ valós, a
		képzetes tengelyen $\tanh(iy) = i \tan y$ tisztán képzetes.}
	\label{fig:tangent-surfaces}
\end{figure}

Az ábrák (\ref{fig:moduli}--\ref{fig:tangent-surfaces}.~ábra) és a tétel
képletei alapján a hat függvény összevetése:
\begin{itemize}
	\item \emph{Zérushelyek.} A $\sin z$ zérushelyei a $k\pi$, a $\cos z$-éi a
	      $\pi/2 + k\pi$ valós számok; a $\sinh z$ zérushelyei a $k\pi i$, a
	      $\cosh z$-éi a $(\pi/2 + k\pi)i$ tisztán képzetes számok
	      $(k \in \mathbb{Z})$. A tangensnek a szinusszal közösek a
	      zérushelyei, és a koszinusz zérushelyein pólusa van; a tangens
	      hiperbolikusz ugyanígy a szinusz és a koszinusz hiperbolikuszból.
	\item \emph{Periodicitás.} A szinusz és a koszinusz a valós tengely
	      irányában periodikus: $\sin(z + 2\pi) = \sin z$. A hiperbolikus
	      függvények a képzetes irányban: $\sinh(z + 2\pi i) = \sinh z$, mert
	      $e^{z + 2\pi i} = e^z$. A tangensek periódusa ennek fele: $\pi$,
	      illetve $\pi i$.
	\item \emph{Növekedés.} A szinusz és a koszinusz a valós tengely mentén
	      korlátos, a képzetes irányban viszont $\frac{1}{2}e^{|y|}$-hez
	      hasonlóan nő; a hiperbolikus függvényeknél éppen fordítva. A
	      tangensek ezzel szemben korlátosak a pólusaiktól távol: a tétel
	      képlete szerint $|\tan(x + iy)| \to 1$, ha $|y| \to \infty$, mert
	      a számláló és a nevező is $\sinh^2 y$-hoz hasonlóan nő; a
	      $|\tanh(x + iy)|$ ugyanígy $1$-hez tart, ha $|x| \to \infty$.
	\item \emph{Elforgatás.} Mivel $|\sinh z| = |\sin(iz)|$,
	      $|\cosh z| = |\cos(iz)|$ és $|\tanh z| = |\tan(iz)|$, és az $i$-vel való szorzás $90^\circ$-os
	      elforgatás, a hiperbolikus függvények képe a trigonometrikusaké
	      $90^\circ$-kal elforgatva.
	\item \emph{Eltolás.} Ahogy a valós számokon, a koszinusz a szinusz
	      eltoltja ($\cos z = \sin(z + \pi/2)$), és ugyanígy a koszinusz
	      hiperbolikusz a szinusz hiperbolikuszé a képzetes irányban.
\end{itemize}

\section{Összefoglalás}

\begin{center}
	\renewcommand{\arraystretch}{2.4}
	\begin{tabular}{@{}lcc@{}}
		\toprule
		                                & Trigonometrikus                    & Hiperbolikus \\
		\midrule
		Görbe                           & $x^2 + y^2 = 1$                    &
		$x^2 - y^2 = 1$                                                                     \\
		Exponenciális alak              &
		$\cos z = \dfrac{e^{iz} + e^{-iz}}{2}$ &
		$\cosh z = \dfrac{e^{z} + e^{-z}}{2}$                                                      \\
		                                & $\sin z = \dfrac{e^{iz} - e^{-iz}}{2i}$ &
		$\sinh z = \dfrac{e^{z} - e^{-z}}{2}$                                                      \\
		Pitagorasz-tétel                & $\cos^2 z + \sin^2 z = 1$          &
		$\cosh^2 z - \sinh^2 z = 1$                                                         \\
		Kapcsolat                       & $\cos z = \cosh(iz)$               &
		$\cosh z = \cos(iz)$                                                                \\
		                                & $i \sin z = \sinh(iz)$             &
		$i \sinh z = \sin(iz)$                                                              \\
		\bottomrule
	\end{tabular}
\end{center}
\end{document}
